\documentclass[ECP,preprint]{ejpecp}

\SHORTTITLE{A variance-scaled Tur\'an inequality at first descent}

\TITLE{A variance-scaled Tur\'an inequality at the first descent of a
  Poisson--binomial mass function%
  \thanks{Generative AI systems assisted with this work; see the disclosure
    and audit status at the end.}}

\AUTHORS{%
  Brett~Reynolds\footnote{Humber College, Toronto, Canada.
    \EMAIL{brett.reynolds@humber.ca}}%
  \orcid{0000-0003-2407-9448}}

\KEYWORDS{Poisson--binomial law; Bernoulli sum; probability mass function;
  log-concavity; Tur\'an inequality; modal index; computer-assisted proof}

\AMSSUBJ{60E15}
\AMSSUBJSECONDARY{60C05; 05A20}

\ABSTRACT{Let $W$ be a finite sum of independent Bernoulli summands with
  probability mass function (pmf) $f=(f_k)$, where
  $f_k=\mathbb{P}(W=k)$, and variance
  $V=\operatorname{Var}(W)\geq1$. If $D$ is the first-descent index of this
  pmf, we prove
  \[
    V\left(1-\frac{f_{D-1}f_{D+1}}{f_D^2}\right)\geq\frac14.
  \]
  This gives a variance-scaled bound for the normalized Tur\'an deficit at
  $D$ and yields $f_{D+r}/f_D\leq\exp(-r/(4V))$ for every support index
  $D+r$ with $r\geq1$.
  An explicit family of finite binomial laws shows that no universal constant
  larger than $1/3$ is possible. The proof derives left and right normalized-mass
  lower bounds about the rightmost modal index from cubic mass inequalities,
  combines them with a lattice maximal-mass bound, and proves the remaining
  scalar inequality directly outside a compact range and by exact Bernstein
  expansions with strictly positive rational coefficients within it.}

\hypersetup{
  pdftitle={A variance-scaled Turan inequality at the first descent of a
    Poisson-binomial mass function},
  pdfauthor={Brett Reynolds},
}

\newcommand{\Pp}{\mathbb{P}}
\newcommand{\Var}{\operatorname{Var}}
\newcommand{\Ber}{\operatorname{Bernoulli}}

\begin{document}

\section{Introduction and main results}\label{sec:introduction}

Let
\[
  W=\sum_{i=1}^n B_i,
  \qquad B_i\sim\Ber(p_i),
\]
where the summands are independent. The law of $W$ is Poisson--binomial;
see \cite{tangtang2023} for a recent survey.

Bernoulli summands with $p_i=0$ may be deleted, whereas those with $p_i=1$
contribute a deterministic shift. Put
$I=\{i:0<p_i<1\}$,
$s=\#\{i:p_i=1\}$, and
$W^\circ=\sum_{i\in I}B_i$. Then $W=s+W^\circ$ and
$\Var(W^\circ)=\Var(W)$. Let $h=(h_k)$ be the pmf of $W^\circ$, where
$h_k=\Pp(W^\circ=k)$. Then $\Pp(W=s+k)=h_k$; hence translation shifts the
first-descent index by $s$ and
preserves the corresponding adjacent mass ratios and normalized Tur\'an
deficits wherever defined.

Assume $\Var(W)\geq1$. Then $I$ is nonempty. We henceforth replace $W$ by
$W^\circ$, relabel $|I|$ as $n$, and assume that every success probability
satisfies $0<p_i<1$. Put
\[
  f_k=\Pp(W=k),\qquad V=\Var(W)=\sum_{i=1}^n p_i(1-p_i),
\]
and call $(f_k)_{k=0}^n$ the probability mass function (pmf) of $W$. Extend
the pmf by $f_{-1}=f_{n+1}=0$. The probability-generating polynomial
$\sum_{k=0}^n f_kx^k$ has only negative real zeros, so Newton's inequalities
make the pmf log-concave \cite[equation~(1)]{pitman1997}.
For $0\leq k\leq n$, define the normalized Tur\'an deficit
\begin{equation}\label{eq:deficit}
  \delta_k
  :=1-\frac{f_{k-1}f_{k+1}}{f_k^2}.
\end{equation}
We have $0\leq\delta_k\leq1$ and $\delta_0=\delta_n=1$.

Let
\begin{equation}\label{eq:first-descent}
  D:=\min\{1\leq k\leq n:f_k<f_{k-1}\}.
\end{equation}
This set is nonempty. Indeed,
\[
  \frac{f_n}{f_{n-1}}
  =\left(\sum_{i=1}^n\frac{1-p_i}{p_i}\right)^{-1}.
\]
If $f_n\geq f_{n-1}$, then the sum on the right is at most one, whereas
\[
  V=\sum_{i=1}^n p_i^2\frac{1-p_i}{p_i}
  <\sum_{i=1}^n\frac{1-p_i}{p_i}\leq1,
\]
a contradiction. By log-concavity, $c:=D-1$ is the rightmost modal index.

Ordinary log-concavity gives only $\delta_D\geq0$. The sharper
$\operatorname{ULC}(n)$ estimate scales with the number of summands, which can
grow while the variance remains fixed. Darroch's localization theorem controls
modal-index location; variance-scaled maximal-mass bounds control the maximal
mass; neither controls the normalized Tur\'an deficit. The cubic mass
inequalities of Hillion and Johnson supply a recurrence for that deficit but
don't compare it with variance
\cite{bobkovmarsigliettimelbourne2022,darroch1964,hillionjohnson2016}.

The variance of the Bernoulli summand $B_i$ is $p_i(1-p_i)$, which tends to
zero as $p_i$ approaches $0$ or $1$. Thus the natural local question is
whether $V\delta_D$ is uniformly bounded away from zero. The answer is yes.

\begin{theorem}\label{thm:main}
  Under the standing notation above, suppose $V\geq1$ and let $D$ be defined
  by \eqref{eq:first-descent}. Then
  \begin{equation}\label{eq:main-bound}
    V\delta_D
    =V\left(1-\frac{f_{D-1}f_{D+1}}{f_D^2}\right)
    \geq\frac14.
  \end{equation}
\end{theorem}

A positive universal bound without a lower variance threshold is impossible.
If $W\sim\Ber(p)$ with $0<p<1/2$, then $D=1$, $\delta_D=1$, and
$V\delta_D=p(1-p)\to0$ as $p\downarrow0$.

Let $c_\star$ denote the infimum of $V\delta_D$ over all finite
Poisson--binomial laws with $V\geq1$, where $D$ and $\delta_D$ are defined
from each law's pmf as above. The theorem and the next proposition give
\begin{equation}\label{eq:constant-bracket}
  \frac14\leq c_\star\leq\frac13.
\end{equation}

\begin{proposition}\label{prop:one-third}
  For every $n\geq5$, there exists a binomial law of variance one whose pmf has
  first-descent index $D=2$ and for which
  \[
    \delta_D=\frac{n+1}{3(n-1)}.
  \]
  This gives $c_\star\leq1/3$.
\end{proposition}

\begin{proof}
  Set
  \[
    p_n=\frac{1-\sqrt{1-4/n}}2,
    \qquad W_n\sim\operatorname{Bin}(n,p_n).
  \]
  Then $np_n(1-p_n)=1$. Write $f=(f_k)$ for the pmf of $W_n$ and put
  $q_k=f_k/f_{k-1}$. Then
  \[
    q_k=\frac{n-k+1}{k}\frac{p_n}{1-p_n}.
  \]
  Since $p_n(1-p_n)=1/n$ and $p_n<1$, we have $p_n>1/n$ and hence
  $q_1>1$. Also $p_n<2/(n+1)$ for $n\geq5$: on $[0,1/2]$ the map
  $x\mapsto x(1-x)$ is increasing, and
  \[
    \frac{2}{n+1}\left(1-\frac{2}{n+1}\right)
    =\frac{2(n-1)}{(n+1)^2}>\frac1n.
  \]
  So $q_2<1$ and $D=2$. Finally,
  \[
    \delta_2=1-\frac{q_3}{q_2}
    =1-\frac{2(n-2)}{3(n-1)}
    =\frac{n+1}{3(n-1)}\longrightarrow\frac13.
  \]
\end{proof}

Proposition~\ref{prop:one-third} supplies a limiting admissible family but no
extremality argument. The lower-bound proof also uses several relaxations. It
first derives explicit left and right normalized-mass lower bounds. When the
scalar $H$ introduced in
Section~\ref{sec:scalar} satisfies $H\geq4$, the proof replaces each right
normalized-mass lower bound by a smaller symmetric normalized-mass lower bound
before applying the general maximal-mass bound. The inequalities defining
these normalized-mass lower bounds and the maximal-mass bound need not be
equalities simultaneously. The argument establishes
\eqref{eq:constant-bracket}, and we make no conjecture about which,
if either, endpoint equals $c_\star$.

Theorem~\ref{thm:main} also bounds the adjacent-ratio drop
$1-f_{D+1}/f_D$; log-concavity then bounds $f_{D+r}/f_D$ at every support
index $D+r>D$.
Complementation followed by translation extends these conclusions to
differences of independent Bernoulli sums.

\begin{corollary}\label{cor:signed}
  Under the hypotheses of Theorem~\ref{thm:main},
  \begin{equation}\label{eq:adjacent-ratio-drop}
    V\left(1-\frac{f_{D+1}}{f_D}\right)\geq\frac14.
  \end{equation}
  Writing $\rho_V=1-(4V)^{-1}$, we also have
  \begin{equation}\label{eq:right-tail-decay}
    \frac{f_{D+r}}{f_D}
    \leq\rho_V^r
    \leq\exp\left(-\frac{r}{4V}\right),
    \qquad 1\leq r\leq n-D.
  \end{equation}
  More generally, let $X$ and $Y$ be finite Bernoulli sums, with all summands
  in the two sums mutually independent, and put $Z=X-Y$. Write $g=(g_j)$ for
  the pmf of $Z$, so $g_j=\Pp(Z=j)$, and assume $V_Z=\Var(Z)\geq1$. Define
  $d:=\min\{j\in\mathbb Z:g_j>0\text{ and }g_j<g_{j-1}\}$. Complementation
  followed by translation to a Bernoulli sum shows that this index is well
  defined. If
  $u=\max\{j:g_j>0\}$ and $\rho_Z=1-(4V_Z)^{-1}$, then
  \[
    V_Z\left(1-\frac{g_{d-1}g_{d+1}}{g_d^2}\right)\geq\frac14,
    \qquad
    V_Z\left(1-\frac{g_{d+1}}{g_d}\right)\geq\frac14.
  \]
  We also have
  \[
    \frac{g_{d+r}}{g_d}
    \leq\rho_Z^r
    \leq\exp\left(-\frac{r}{4V_Z}\right),
    \qquad 1\leq r\leq u-d.
  \]
\end{corollary}

\begin{proof}
  Write $\eta_-=f_D/f_{D-1}<1$ and $\eta_+=f_{D+1}/f_D$. Then
  \[
    1-\eta_+\geq1-\frac{\eta_+}{\eta_-}=\delta_D,
  \]
  proving \eqref{eq:adjacent-ratio-drop}. Log-concavity makes the adjacent
  mass ratios nonincreasing. Hence, for $1\leq j\leq r$,
  \[
    \frac{f_{D+j}}{f_{D+j-1}}
    \leq\frac{f_{D+1}}{f_D}\leq\rho_V.
  \]
  Multiplication gives the first inequality in
  \eqref{eq:right-tail-decay}, and $1-x\leq e^{-x}$ gives the second.

  Write $Y=\sum_{j=1}^mY_j$. Then
  \[
    Z+m=X+\sum_{j=1}^m(1-Y_j)
  \]
  is a finite Bernoulli sum. Translation preserves the variance and the
  corresponding adjacent mass ratios and normalized Tur\'an deficits while
  shifting the first-descent index, so the same argument gives all three
  signed estimates.
\end{proof}

Equation~\eqref{eq:right-tail-decay} gives a variance-scaled geometric upper
bound for the normalized mass $f_{D+r}/f_D$ at every support index $D+r$ with
$1\leq r\leq n-D$. The standard
$\operatorname{ULC}(n)$ estimate for Poisson--binomial pmfs alone
doesn't recover this conclusion:
\begin{equation}\label{eq:ulc}
  \delta_k\geq
  \frac{n+1}{(k+1)(n-k+1)};
\end{equation}
see \cite[Definition~3.11 and equation~(41)]{hillionjohnson2016} and, for a
recent treatment via relative log-concavity,
\cite[the discussion following Definition~1.3]{marsigliettimelbourne2026}.
This bound is indexed by the number of summands rather than by the variance
$V$ of their sum.

For example, take $m$ success probabilities equal to $\varepsilon_m$ and $m$
equal to $1-\varepsilon_m$, where
$2m\varepsilon_m(1-\varepsilon_m)=1$. The resulting pmf is symmetric and
strictly log-concave. The law has $V=1$, and the pmf has first-descent index
$D=m+1$, but the right-hand side of
\eqref{eq:ulc} at $D$ is
\[
  \frac{2m+1}{m(m+2)}\longrightarrow0.
\]
Even at fixed variance, the $\operatorname{ULC}(n)$ lower bound at the
first-descent index tends to zero.

Existing quantitative results control modal-index location, adjacent mass
ratios, or individual masses rather than the variance-scaled normalized
Tur\'an deficit at the first-descent index. Darroch shows that a modal index
lies within one unit of the mean \cite{darroch1964}. Pitman bounds adjacent
mass ratios through exponentially tilted real-rooted laws
\cite[equations~(20)--(21)]{pitman1997}, and D\"umbgen and Wellner prove strict
decrease of the index-weighted mass ratios
$(k+1)f_{k+1}/f_k$ \cite[Proposition~2]{duembgenwellner2020}.

Gnedin extends the mean--mode rule to extended Bernoulli sums
\cite[Theorem~2]{gnedin2024}. A three-mass peak-skewness statistic measures
the Bernoulli perturbation needed to change the leading mode
\cite[Sections~4.2 and~7]{gnedin2024}, rather than the normalized Tur\'an
deficit at the first-descent index considered here.

Johnson's $c$-log-concavity bound also controls $\delta_k$, but its scale
depends on the success probabilities:
\[
  \delta_k\geq \frac{f_{k+1}}{f_k}
  \left(\sum_i\frac{p_i}{1-p_i}\right)^{-1};
\]
see \cite[Definition~1.2, Lemma~5.1, and
Example~5.2(2)]{johnson2017discrete}. In the variance-one symmetric family
above, the reciprocal of the odds sum tends to zero. At the first-descent
index, the adjacent mass ratio is at most one, so the right-hand side of
Johnson's bound also tends to zero.

The uniform maximal-mass bound of Baillon, Cominetti, and Vaisman and the
lattice maximal-mass bound of Bobkov, Marsiglietti, and Melbourne control one
mass rather than the normalized Tur\'an deficit in \eqref{eq:deficit}
\cite{bailloncominettivaisman2016,bobkovmarsigliettimelbourne2022}. We aren't
aware of an earlier universal constant $c>0$ for which
$V\delta_D\geq c$ at the first-descent index of the pmf of every finite
Poisson--binomial law with $V\geq1$.

For an integer-valued log-concave random variable $X$ with maximal mass $M$,
Aravinda's sharp min-entropy/variance inequality gives
\cite[Theorem~1.4]{aravinda2024}
\[
M^{-2}\geq1+\Var(X).
\]
This bounds the maximal mass from above. Its direction is opposite to the input needed here:
\eqref{eq:maximal-mass-bound} gives
$M\geq(1+12V)^{-1/2}$, a lower bound on the maximal mass. Aravinda's result
complements that step and reinforces the distinction between variance control
of one atom and variance control of a local normalized Tur\'an deficit.

The proof uses two external results. The cubic mass inequalities of Hillion
and Johnson yield a recurrence; iteration yields explicit left and right
normalized-mass lower bounds about the rightmost modal index. The
pairwise variance identity turns those bounds into a variance estimate;
combining that estimate with the maximal-mass bound of Bobkov, Marsiglietti,
and Melbourne reduces the theorem to a scalar inequality. A direct symbolic
argument proves that inequality for $H\geq16$; exact Bernstein expansions with
strictly positive rational coefficients prove it for $3<H\leq16$, without
floating-point sampling.

\section{The recurrence for normalized Tur\'an deficits}
\label{sec:propagation}

For $1\leq k\leq n$, put
\[
  q_k=\frac{f_k}{f_{k-1}}.
\]
Log-concavity says that $(q_k)$ is nonincreasing, and for
$1\leq k\leq n-1$,
\begin{equation}\label{eq:delta-ratio}
  \delta_k=1-\frac{q_{k+1}}{q_k}.
\end{equation}

Hillion and Johnson's Theorem~A.2 and Corollary~A.3 state, for every
$k\in\mathbb Z$, the cubic mass inequalities
\cite[(78)--(79)]{hillionjohnson2016}
\begin{align}
  (f_{k-1}^2-f_{k-2}f_k)f_{k+1}
  &\leq f_{k-1}(f_k^2-f_{k-1}f_{k+1}),
  \label{eq:HJ-left}\\
  (f_{k+1}^2-f_kf_{k+2})f_{k-1}
  &\leq f_{k+1}(f_k^2-f_{k-1}f_{k+1}).
  \label{eq:HJ-right}
\end{align}
Their all-integer statement understands the pmf as zero beyond its support,
so these inequalities include the cases adjacent to a support boundary. The
identities
\[
\begin{aligned}
  f_{k-1}^2-f_{k-2}f_k&=f_{k-1}^2\delta_{k-1},\\
  f_k^2-f_{k-1}f_{k+1}&=f_k^2\delta_k,\\
  f_{k+1}^2-f_kf_{k+2}&=f_{k+1}^2\delta_{k+1}
\end{aligned}
\]
show the normalization explicitly. For $1\leq k\leq n-1$, dividing
\eqref{eq:HJ-left} by $f_{k-1}f_k^2$ and \eqref{eq:HJ-right} by
$f_{k+1}f_k^2$, respectively, gives the recurrence for normalized Tur\'an
deficits
\begin{equation}\label{eq:deficit-recurrence}
  \delta_{k-1}(1-\delta_k)\leq\delta_k,
  \qquad
  \delta_{k+1}(1-\delta_k)\leq\delta_k.
\end{equation}
In particular, this includes relations adjacent to the boundary, such as
$\delta_n(1-\delta_{n-1})\leq\delta_{n-1}$.

Set $\delta=\delta_D$. The boundary values show that $\delta>0$: if
$\delta=0$, \eqref{eq:deficit-recurrence} forces both adjacent values of
$\delta_k$ to vanish; repeated application then forces every $\delta_k$ to
vanish, contradicting
$\delta_0=\delta_n=1$. If $\delta\geq1/4$, Theorem~\ref{thm:main} follows
from $V\geq1$, so it remains to consider
\begin{equation}\label{eq:small-delta}
  0<\delta<\frac14.
\end{equation}
In this case $D$ lies strictly inside the support: $D\geq1$ by definition,
whereas $\delta_D<1=\delta_n$ forces $D<n$.

Iterating the map $x\mapsto x/(1-x)$ in \eqref{eq:deficit-recurrence} gives,
for either sign and every integer $r\geq1$,
\begin{equation}\label{eq:deficit-bound}
  \delta_{D\pm r}\leq\frac{\delta}{1-r\delta}
  \qquad\text{provided }D\pm r\in\{0,\ldots,n\}
  \text{ and }(r+1)\delta<1.
\end{equation}
The same bound excludes an endpoint index of the support in this range:
otherwise the left-hand side would equal one, while the right-hand side is
strictly less than one.

Define
\begin{equation}\label{eq:K-a}
  K=\max\{r\in\mathbb Z_{\geq1}:(r+1)\delta<1\},
  \qquad
  a=\frac{1-2\delta}{1-\delta}.
\end{equation}
The endpoint exclusion gives $D>K$ and $n-D>K$, so every index used below
lies in the support and its corresponding mass is strictly positive. Recall
that
$c=D-1$ is the rightmost modal index, and write $p=f_c=\max_k f_k$ for the
maximal mass. Since $q_{D-1}\geq1>q_D$, the first inequality in
\eqref{eq:deficit-recurrence} gives
\[
  \delta_{D-1}\leq\frac{\delta}{1-\delta},
  \qquad
  q_D=q_{D-1}(1-\delta_{D-1})\geq a.
\]

For $0\leq j\leq K-1$, \eqref{eq:deficit-bound} gives
\[
\begin{aligned}
  q_{D+j}
  &=q_D\prod_{s=0}^{j-1}(1-\delta_{D+s})\\
  &\geq a\prod_{s=0}^{j-1}
    \frac{1-(s+1)\delta}{1-s\delta}
   =a(1-j\delta).
\end{aligned}
\]
For $1\leq r\leq K$, we obtain the right normalized-mass lower bound
\begin{equation}\label{eq:right-mass-bound}
  R_r:=a^r\prod_{j=1}^{r-1}(1-j\delta),
  \qquad
  \frac{f_{c+r}}p=\prod_{j=0}^{r-1}q_{D+j}\geq R_r.
\end{equation}
For the left normalized-mass lower bounds, let $1\leq j\leq K$. Then
\[
  q_D=q_{D-j}\prod_{s=1}^j(1-\delta_{D-s})<1,
\]
and hence
\[
\begin{aligned}
  \frac1{q_{D-j}}
  &>\prod_{s=1}^j(1-\delta_{D-s})\\
  &\geq\prod_{s=1}^j
    \frac{1-(s+1)\delta}{1-s\delta}
   =\frac{1-(j+1)\delta}{1-\delta}.
\end{aligned}
\]
For $1\leq r\leq K$, we obtain the left normalized-mass lower bound
\begin{equation}\label{eq:left-mass-bound}
  L_r:=(1-\delta)^{-r}
  \prod_{j=2}^{r+1}(1-j\delta),
  \qquad
  \frac{f_{c-r}}p=\prod_{j=1}^r\frac1{q_{D-j}}\geq L_r.
\end{equation}

\section{Variance bounds and scalar reduction}
\label{sec:variance}

Define auxiliary weights by $w_0=1$, $w_r=R_r$, and $w_{-r}=L_r$ for
$1\leq r\leq K$, and put
\begin{equation}\label{eq:A-def}
  A(\delta)=\frac12\sum_{i,j=-K}^K w_iw_j(i-j)^2.
\end{equation}
The pairwise variance identity and the preceding normalized-mass lower bounds
give
\begin{equation}\label{eq:pairwise-variance}
  V=\frac12\sum_{i,j=0}^n f_if_j(i-j)^2
  \geq p^2A(\delta).
\end{equation}

We also use the lattice maximal-mass bound of Bobkov, Marsiglietti, and
Melbourne
\cite[Corollary~3.2]{bobkovmarsigliettimelbourne2022}:
\begin{equation}\label{eq:maximal-mass-bound}
  V\geq\frac{p^{-2}-1}{12}.
\end{equation}
A short smoothing argument proves \eqref{eq:maximal-mass-bound}. Let $U$ be
independent of $W$ and uniform on the unit interval centered at zero. The
density of $W+U$ is bounded by $p$, and
$\Var(W+U)=V+1/12$. Let $\varphi$ be any probability density bounded by $p$,
and let $x_0\in\mathbb R$. Its second moment about $x_0$ is minimized by
filling the interval of length $1/p$ centered at $x_0$ at density $p$.
Moving mass toward $x_0$ without exceeding the density bound lowers the
integral until this density is reached. Hence
\[
  \int_{\mathbb R}(x-x_0)^2\varphi(x)\,dx
  \geq p\int_{-1/(2p)}^{1/(2p)}t^2\,dt
  =\frac1{12p^2}.
\]
Taking $x_0=\mathbb E(W+U)$ yields $V+1/12\geq1/(12p^2)$, proving
\eqref{eq:maximal-mass-bound}.

It remains to prove one scalar estimate.

\begin{proposition}\label{prop:scalar}
  For $0<\delta<1/4$, the quantity in \eqref{eq:A-def} satisfies
  \begin{equation}\label{eq:scalar-target}
    A(\delta)\geq\frac{3+\delta}{4\delta^2}.
  \end{equation}
\end{proposition}

\begin{proof}[Proof of Theorem~\ref{thm:main} from
  Proposition~\ref{prop:scalar}]
  Suppose that $V<1/(4\delta)$. Inequality
  \eqref{eq:maximal-mass-bound} then gives
  \[
    p^2>\frac{\delta}{3+\delta}.
  \]
  On the other hand, \eqref{eq:pairwise-variance} and
  \eqref{eq:scalar-target} give
  \[
    p^2<\frac1{4\delta A(\delta)}
    \leq\frac{\delta}{3+\delta},
  \]
  a contradiction.
\end{proof}

\section{The scalar inequality}
\label{sec:scalar}

Set
\begin{equation}\label{eq:H-def}
  H=\frac{1-\delta}{\delta}=\frac1\delta-1>3.
\end{equation}
Under \eqref{eq:H-def}, the left normalized-mass lower bounds are
\begin{equation}\label{eq:b-def}
  L_r=b_r:=\prod_{s=1}^r\left(1-\frac{s}{H}\right).
\end{equation}
For $1\leq r\leq K$, the right normalized-mass lower bounds satisfy
$R_r\geq L_r=b_r$. Indeed,
if $C_r=R_r/L_r$, then $C_1=1$
and, for $1\leq r<K$,
\begin{equation}\label{eq:C-ratio}
  \frac{C_{r+1}}{C_r}
  =\frac{(H-1)(H+1-r)}{(H+1)(H-r-1)}
  =1+\frac{2r}{(H+1)(H-r-1)}\geq1.
\end{equation}

Since every auxiliary weight is nonnegative, the pairwise form in
\eqref{eq:A-def} is nondecreasing in each auxiliary weight. With $K$ fixed,
\[
  \frac{\partial A}{\partial w_\ell}
  =\sum_{j=-K}^K w_j(\ell-j)^2\geq0.
\]
Replacing every $R_r$ by $b_r$ doesn't increase the pairwise form; denote the
resulting value by $A_{\mathrm{sym}}$. For the symmetric auxiliary weights
$w_0=1$ and $w_{\pm r}=b_r$, put
\begin{equation}\label{eq:S-T}
  S=1+2\sum_{r=1}^K b_r,
  \qquad
  T=2\sum_{r=1}^K r^2b_r.
\end{equation}
The pairwise identity gives
\[
  A_{\mathrm{sym}}
  =\left(\sum_{r=-K}^K w_r\right)
   \left(\sum_{r=-K}^K r^2w_r\right)
   -\left(\sum_{r=-K}^K rw_r\right)^2
  =ST,
\]
because the final sum vanishes by symmetry. Hence
$A(\delta)\geq A_{\mathrm{sym}}=ST$. Under \eqref{eq:H-def}, the right-hand
side of \eqref{eq:scalar-target} is
\begin{equation}\label{eq:Q-def}
  Q(H)=\frac{(3H+4)(H+1)}4.
\end{equation}
For $H\geq4$, it remains to prove $ST\geq Q(H)$. On $3<H<4$, we instead
retain the exact asymmetric normalized-mass lower bounds.

\subsection{The compact range}

On $3<H\leq4$, use the exact normalized-mass lower bounds
$L_1,L_2,L_3,R_1,R_2,R_3$. Direct simplification gives
\begin{equation}\label{eq:asymmetric-P}
  A(\delta)-Q(H)=\frac{P(H)}{4H^5(H+1)^3},
\end{equation}
where
\begin{align*}
  P(H)={}&-3H^{10}-16H^9+750H^8-3676H^7+6613H^6
  -5460H^5\\
  &+800H^4+4696H^3-7176H^2+4272H-912.
\end{align*}
After substituting $H=3+t$, the degree-ten Bernstein expansion of $P$ on
$0\leq t\leq1$ has eleven strictly positive rational coefficients.

For $m=4,\ldots,15$ and $H\in[m,m+1]$, define
\begin{equation}\label{eq:compact-ST}
  S_m=1+2\sum_{r=1}^m b_r,
  \qquad
  T_m=2\sum_{r=1}^m r^2b_r,
\end{equation}
and
\begin{equation}\label{eq:compact-P}
  P_m(H)=4H^{2m}\bigl(S_mT_m-Q(H)\bigr).
\end{equation}
For $m<H\leq m+1$, we have $K=m$. At $H=m$, we have $K=m-1$ and
$b_m=0$. Hence $S_m=S$ and $T_m=T$ throughout the closed cell
$[m,m+1]$. Each $P_m$ is an integer polynomial of degree $2m+2$. After
substituting $H=m+t$, all its full-degree Bernstein coefficients are strictly
positive.

If $G(t)=\sum_{j=0}^d a_jt^j$ has degree $d$, the Bernstein conversion used
here is
\begin{equation}\label{eq:bernstein-conversion}
  \beta_i=\sum_{j=0}^i a_j
  \frac{\binom{i}{j}}{\binom{d}{j}},
  \qquad
  G(t)=\sum_{i=0}^d \beta_i\binom di t^i(1-t)^{d-i}.
\end{equation}
All polynomial identities and Bernstein coefficients in this step are exact
and rational.
The Bernstein expansions of $P(3+t)$ and the twelve polynomials $P_m(m+t)$,
$4\leq m\leq15$, have, respectively, $11$ and $264$ strictly positive
coefficients. Thus all $275$ certificate coefficients are strictly positive.

\begin{table}[t]
  \centering
  \caption{Compact-range certificate summary. The minimum is taken over the
  Bernstein coefficients of the stated cleared numerator; the final column
  gives the first sixteen hexadecimal characters of its coefficient-vector
  SHA-256 digest. Full coefficients and digests are in the supplement.}
  \label{tab:certificate-summary}
  \scriptsize
  \setlength{\tabcolsep}{4pt}
  \begin{tabular}{lrrrl}
    \hline
    Cell & Degree & Count & Minimum coefficient & SHA-256 prefix \\
    \hline
    $[3,4]$ asym. & 10 & 11 & $22272$ & \texttt{e1e452c9c93c30cd} \\
    $[4,5]$ & 10 & 11 & $332800$ & \texttt{441f603ac3f60919} \\
    $[5,6]$ & 12 & 13 & $476945150$ & \texttt{80a652bad6b13e69} \\
    $[6,7]$ & 14 & 15 & $252843503616$ & \texttt{e3ad5553e47317b2} \\
    $[7,8]$ & 16 & 17 & $141578177942152$ & \texttt{6c774572e048423c} \\
    $[8,9]$ & 18 & 19 & $92316620768411648$ & \texttt{e9b192d522d49b1e} \\
    $[9,10]$ & 20 & 21 & $71284663153149460650$ & \texttt{876f43f48d0c3c1e} \\
    $[10,11]$ & 22 & 23 & $65053547560474378240000$ & \texttt{b741906966aed963} \\
    $[11,12]$ & 24 & 25 & $69645362929101036447979796$ & \texttt{6c318dad110490486} \\
    $[12,13]$ & 26 & 27 & $86706679760909016273411637248$ & \texttt{643c6c7cf85a5535} \\
    $[13,14]$ & 28 & 29 & $124437022735783915561664544889462$ & \texttt{64f915da71a6e95d} \\
    $[14,15]$ & 30 & 31 & $204172536352322626071425358548172800$ & \texttt{e31285e840aa69d9} \\
    $[15,16]$ & 32 & 33 & $380093996939768323066638847260414000000$ & \texttt{0850f4a393088df0} \\
    \hline
  \end{tabular}
\end{table}

The generator program reconstructs every identity before checking the signs.
This reconstruction provides an internal consistency check.
The accompanying full certificate data file records each exact power-basis
numerator and all $275$ reduced rational Bernstein coefficients. A separately
implemented standard-library checker supplies a computational cross-check: it
reconstructs the identities and Bernstein expansions for the thirteen compact
scalar cells without importing the generator program. Because the Bernstein
basis functions are nonnegative on $[0,1]$, the verified coefficient signs prove
Proposition~\ref{prop:scalar} for $3<H\leq16$.

\subsection{The analytic range}

Thirteen exact scalar-cell checks cover the compact range; one initial
triangular cell and a uniform cell family cover the analytic range.

Suppose $H\geq16$, and choose an integer $J\geq5$ such that
\begin{equation}\label{eq:triangular-cell}
  \frac{J(J+1)}2\leq H\leq\frac{(J+1)(J+2)}2.
\end{equation}
Either choice is valid at a shared endpoint. Since $\delta=(H+1)^{-1}$,
\[
  K=\max\{r\in\mathbb Z_{\geq1}:r<H\}.
\]
For $J\geq5$, $J<J(J+1)/2\leq H$, so $J\leq K$. For
$1\leq r\leq J$, the elementary product bound gives
\begin{equation}\label{eq:bonferroni}
  b_r\geq\lambda_r:=1-\frac{r(r+1)}{2H}\geq0.
\end{equation}
Put
\begin{align}
  S_0&=1+2\sum_{r=1}^J \lambda_r
  =2J+1-\frac{J(J+1)(J+2)}{3H},
  \label{eq:S0}\\
  T_0&=2\sum_{r=1}^Jr^2\lambda_r
  =\frac{J(J+1)(2J+1)}3-\frac{M_4+M_3}{H},
  \label{eq:T0}
\end{align}
where
\[
  M_3=\frac{J^2(J+1)^2}{4},
  \qquad
  M_4=\frac{J(J+1)(2J+1)(3J^2+3J-1)}{30}.
\]
Since $J\leq K$, all omitted auxiliary weights are nonnegative and
$b_r\geq\lambda_r$ for $r\leq J$. Hence $S\geq S_0$, $T\geq T_0$, and
$A(\delta)\geq ST\geq S_0T_0$. It remains to prove $S_0T_0\geq Q(H)$.

Let $N_J(H)=H^2(S_0T_0-Q(H))$. For $J=5$, only
$16\leq H\leq21$ is needed. After writing $H=16+5t$, $0\leq t\leq1$, the
degree-four Bernstein coefficients are
\[
  2360,\qquad 7500,\qquad \frac{25055}{2},\qquad
  \frac{254205}{16},\qquad \frac{31115}{2}.
\]

For $J\geq6$, write
\[
  H=\frac{J(J+1)}2+(J+1)t,
  \qquad J=u+6,
  \qquad 0\leq t\leq1.
\]
After removing strictly positive multipliers, the five degree-four Bernstein
coefficients in $t$ of
\[
  N_J\left(\frac{J(J+1)}2+(J+1)t\right)
\]
are the following polynomials in $u$, in order:
\begin{align*}
  &121u^4+2474u^3+17431u^2+46014u+25400,\\
  &121u^5+3442u^4+37967u^3+199405u^2+480475u+384510,\\
  &121u^6+4410u^5+65863u^4+512764u^3
    +2172926u^2+4668316u+3829440,\\
  &121u^5+3684u^4+43735u^3+249961u^2+671859u+644400,\\
  &121u^4+2958u^3+25579u^2+88782u+91440.
\end{align*}
After factoring out the common strictly positive scalar $1/2880$, the removed
multipliers are, respectively,
\[
  J^2(J+1)^2,\quad J(J+1)^2,\quad (J+1)^2,\quad
  (J+1)^2(J+2),\quad (J+1)^2(J+2)^2.
\]
All five multipliers are strictly positive for $J\geq6$, and every coefficient
of the five polynomials above is strictly positive for $u\geq0$. Thus
$N_J(H)>0$ and $S_0T_0>Q(H)$ throughout the analytic range. Together with the
compact range Bernstein verification, this proves Proposition~\ref{prop:scalar}.

\par\bigskip
\begin{supplement}
\stitle{Code and exact certificate data}
\sdescription{The supplementary archive comprises the generator program, an
independently implemented checker, the compact summary certificate, the full
certificate, the MIT license, and a manifest giving the reference environment,
exact replay commands, per-cell and aggregate digests, and archive status. The
full certificate records the exact power-basis numerators and all $275$
reduced rational Bernstein coefficients for the thirteen compact scalar cells.
Using only the Python standard library and without importing the generator,
the checker reconstructs the identities and Bernstein expansions, verifies
strict positivity, and recomputes the Bernstein-vector digest for each cell,
each cell-payload digest, and the overall payload digest. The programs verify
only the thirteen compact scalar cells. Neither program machine-verifies
Theorem~\ref{thm:main}. The archive also contains the complete conditional Lean
project cited in the acknowledgement. That project formalizes the recurrence
propagation, endpoint exclusion, first-crossing algebra, and raw-drop
corollary, taking the normalized Hillion--Johnson recurrence as a hypothesis.}
\end{supplement}

\medskip
\begin{acks}[AI use and author responsibility]
OpenAI's GPT-5.6 model, accessed through Codex, supplied substantial proof-search
assistance. Other sessions with Claude (Anthropic), ChatGPT and Codex (both
OpenAI), and Gemini (Google) assisted with computational exploration,
exact-verification code, proof critique, formalization support, literature
searching, drafting, and revision. Aristotle (Harmonic) completed bounded proof
obligations in Lean for a conditional formalization; that formalization assumes
the normalized Hillion--Johnson recurrence and doesn't formally prove
Theorem~\ref{thm:main}. Its complete Lean project is included in the supplement.
The author remains responsible for the mathematics, sources, wording, and
released files. Fresh generator output matched both certificate files byte for
byte, and the independently implemented checker replayed the full certificate
in the reference environment; the exact commands and digests are listed in the
supplement.
\end{acks}

\clearpage
\bibliographystyle{amsplain}
\bibliography{references}

\end{document}
